\begin{exsol}{23}{\MR{} Food for Thought~2.5}{52}{Exam}{B3}
\asked Assess how specific a cable-television company's investment in wiring a community is,
and predict the arrangements the company would seek with local communities over the pricing and
taxation of its service.

\method Transaction attributes, asset specificity first (\MR{} Ch.~2, pp.~30--31, and summary,
p.~49); imperfect commitment and opportunistic renegotiation (p.~30); specific and cospecialized
assets, the hold-up problem and the contractual remedies observed in coal supply (\MR{} Ch.~5,
``Investments and Specific Assets'', pp.~134--139); quasi-rents (\MR{} Ch.~8, p.~269);
investments in bargaining position (\MR{} p.~149).

\solution
\textbf{(a) How specific is the investment? Very.} \MR{} measure specificity, when value
maximization applies, as ``the percentage of investment value that is lost when the asset is
used outside the specific setting or relationship'' (p.~135). For a cable network:
\begin{itemize}
  \item \emph{Site specificity.} Cable strung on poles or buried along a town's streets to its
  households cannot be moved to another town. The bulk of the cost is installation (trenching,
  pole attachments, drops to each house, engineering), which has no salvage value; only the
  head-end electronics and some equipment are redeployable.
  \item \emph{Dedicated to one counterparty.} The network serves only that community, and the
  firm's access to streets, poles and subscribers runs through a franchise that the local
  government grants, renews and can regulate and tax.
  \item \emph{Durable.} The cable lasts many years, so the investment is recovered only over a
  long horizon, during which conditions (technology, costs, local politics) change.
\end{itemize}
So most of the investment would be lost if the relationship with the community ended: the
specificity is close to $100\%$ for the installed plant. The specificity is \emph{two-sided}, though
not symmetric: once one network is in place, duplicating it is wasteful (cable distribution is a
natural monopoly), so the community also has few alternatives and the firm could exploit its
local monopoly. Ex ante, several firms can compete for the franchise; ex post, firm and community
are locked in a bilateral relationship.

\textbf{(b) The hazard: hold-up.} After the cable is sunk, the community can change the terms:
cap subscriber prices, raise franchise fees or taxes, impose new service obligations, or threaten
not to renew. The firm keeps operating as long as revenue covers its \emph{avoidable} costs, so
the community can appropriate the firm's quasi-rent, the excess of revenue over what the firm
needs to stay rather than exit (\MR{} p.~269). A firm that foresees this underinvests or does not
enter at all. A minimal example: the network costs $K = 100$ (fully specific) and generates
operating returns worth $R = 150$ in present value. Building is efficient ($150 > 100$). If,
after the investment, the returns are split $50/50$ by renegotiation with the community, the firm
gets $75 < 100$ and does not build: $50$ of value is lost although nobody behaves irrationally.
The danger runs both ways: a firm with a captive market can also hold up subscribers with high
prices or poor service. This is \MR's mine-mouth plant problem (pp.~135--137), with the community
in the role of the utility.

A community is also a government: it can change taxes and regulations unilaterally, and its
promises are hard to enforce against it. Imperfect commitment (\MR{} p.~30) is therefore
particularly severe here.

\textbf{(c) Arrangements to expect.} Because the investment is large and specific, the parties
``seek to organize in ways that safeguard those investments'' (\MR{} p.~49). Before investing,
the company would want a long-term, detailed franchise contract with the community, of the kind
Joskow found for coal (\MR{} p.~137):
\begin{enumerate}[label=(\roman*)]
  \item \emph{Long duration}, enough to recover the investment, with renewal procedures that
  presume renewal for satisfactory performance, or compensation for the network (e.g.\ purchase
  at an agreed valuation) if the franchise is not renewed or is given to a rival.
  \item \emph{Pricing rules fixed in advance}: an initial rate schedule with escalator clauses
  tied to cost or price indices, or a price-cap formula, rather than rates set at the community's
  discretion after the cable is laid.
  \item \emph{Taxation fixed by formula}: a franchise fee set as a stated percentage of gross
  revenue, with a ceiling, instead of discretionary taxes that could be raised later.
  \item \emph{Exclusivity or protection against overbuilding} for a period, so that the
  community cannot extract terms by threatening to license a competitor once the firm has sunk
  its cost.
  \item \emph{Obligations on the firm}, protecting the community against the monopoly side of
  the lock-in: build-out to all neighbourhoods, service-quality standards, regulated rates for a
  basic tier, public-access channels.
  \item \emph{Dispute resolution} by arbitration or an outside regulator, and where possible
  commitments at a higher level of government that limit what local authorities can change
  later.
\end{enumerate}
The alternative remedy is the one \MR{} report for GM and Fisher Body (p.~137): unified ownership.
Here that means municipal ownership of the network, with the operation contracted out, which
removes the firm's exposure but brings its own incentive problems.

\textbf{(d) Investments in bargaining position.} Since the terms are renegotiated at each renewal,
both sides have an incentive to spend resources that shift the division without creating value
(\MR{} p.~149): the community may cultivate a credible threat to license a second operator or
build a municipal system; the firm may lobby, litigate, or design its network to raise the
community's switching costs. Such spending is pure waste in total-value terms, and good contracts
reduce it by fixing the renewal terms in advance.

\interp This item connects the Ch.~2 vocabulary (asset specificity, imperfect commitment) to the
hold-up logic developed later in the course. Under the Coase benchmark the efficient network would
be built whatever the bargaining power; with specific assets and incomplete, imperfectly enforceable
contracts, the threat of ex-post renegotiation distorts the ex-ante investment. The transaction
cost approach predicts the governance structure that maximizes total value net of those costs:
long-term contracts with pricing formulas, protected tax treatment and renewal guarantees, or
integration. (For a classic study of cable franchise bidding along these lines, see Williamson,
``Franchise Bidding for Natural Monopolies---in General and with Respect to CATV'', \emph{Bell
Journal of Economics}, 1976.)

\verify Internal consistency: the efficient decision compares $R = 150$ with $K = 100$; the firm
compares its post-renegotiation share $R/2 = 75$ with $K$. Any ex-ante agreement guaranteeing the
firm at least $100$ out of $150$ (e.g.\ a fee formula capping the community's take at $50$)
restores the investment, consistent with the remedies in (c) (\texttt{check\_C.py}).

\answer{The investment is highly specific (site-specific, durable, dedicated to one community;
most of its value is lost outside the relationship), which exposes the company to hold-up through
price controls, taxes or non-renewal once the cable is sunk. Expect long-term franchise agreements
signed before investment: long duration and renewal or buy-out guarantees, rate formulas with
escalators or price caps, franchise fees fixed as a capped share of revenue, exclusivity, and
service obligations protecting the community; or, as the alternative, municipal ownership.}
\end{exsol}
